% GENERATED FILE. Edit canonical lessons, then run npm run generate:books.
% canonical-source-hash: fnv1a64:f82b470a169a0dfe
% canonical-lessons: TA-C315-nataiperu, TA-C315-erpatu, TA-C315-nikalvu, TA-C315-tarpotu, TA-C315-eppotavatu

\chapter{Events Take Place}
\label{ch:ta-events-take-place}
\hlchaptermodality{315}

\begin{chapteropening}
\textbf{By the end of this chapter:} \emph{I can say when and how events took place.}

Complete the last lesson of chapter 315: I can say when and how events took place.
\end{chapteropening}

\section[\texorpdfstring{\ta{நடைபெறு} (naṭaipeṟu)}{நடைபெறு (naṭaipeṟu)}]{\ta{நடைபெறு} (naṭaipeṟu) — to take place}

\label{lesson:TA-C315-nataiperu}



\begin{quote}
Before the new one: say the Tamil for by then, then the Tamil for following this.
\end{quote}

\subsection*{You'll want to know: \ta{நடைபெறு}}
\textbf{\ta{நடைபெறு}} — \emph{naṭaipeṟu} — \textquotedblleft{}to take place\textquotedblright{}.

\begin{quote}
\ta{விழா} \ta{நேற்று} \ta{பள்ளியில்} \ta{நடைபெற்றது}.
\end{quote}

\begin{quote}
\emph{viḻā nēṟṟu paḷḷiyil naṭaipeṟṟatu.}
\end{quote}

\begin{quote}
The ceremony took place at the school yesterday.
\end{quote}

Past (it …): \ta{நடைபெற்றது}, \emph{naṭaipeṟṟatu}.

A formal way to say \ta{நட}, to happen, for planned events.

\begin{cousinweb}
\textbf{\ta{நடைபெறு}} (\emph{naṭaipeṟu}) is a Tamil compound built from \ta{நடை}, \textquotedblleft{}walk, course, conduct\textquotedblright{} (DEDR 3582), and \ta{பெறு}, \textquotedblleft{}to get, obtain\textquotedblright{} (DEDR 4422): to get its course.
\end{cousinweb}

\begin{grammarlens}[title={What you've built}]
One more word for telling what happened.
\end{grammarlens}

\subsection*{Your turn}
\begin{itemize}
\raggedright
  \item \emph{Say it:} \emph{naṭaipeṟu}
  \item \emph{Say it:} \emph{naṭaipeṟu}, once more
  \item \emph{Say it:} \emph{viḻā nēṟṟu paḷḷiyil naṭaipeṟṟatu.} — The ceremony took place at the school yesterday.
  \item \emph{Recall:} say the Tamil for by then, then the Tamil for following this, then say \emph{naṭaipeṟu} again
\end{itemize}

\begin{tcolorbox}[breakable,colback=teal!4,colframe=teal!35!black,title={Before you move on}]
What does \ta{நடைபெறு} mean? (\textquotedblleft{}To take place\textquotedblright{}.) Say it once more.
\end{tcolorbox}
\section[\texorpdfstring{\ta{ஏற்படு} (ēṟpaṭu)}{ஏற்படு (ēṟpaṭu)}]{\ta{ஏற்படு} (ēṟpaṭu) — to arise; to come about}

\label{lesson:TA-C315-erpatu}



\begin{quote}
Before the new one: say the Tamil for following this, then the Tamil for to take place.
\end{quote}

\subsection*{You'll want to know: \ta{ஏற்படு}}
\textbf{\ta{ஏற்படு}} — \emph{ēṟpaṭu} — \textquotedblleft{}to arise; to come about\textquotedblright{}.

\begin{quote}
\ta{சாலையில்} \ta{பெரிய} \ta{நெரிசல்} \ta{ஏற்பட்டது}.
\end{quote}

\begin{quote}
\emph{cālaiyil periya nerical ēṟpaṭṭatu.}
\end{quote}

\begin{quote}
A big traffic jam arose on the road.
\end{quote}

Past (it …): \ta{ஏற்பட்டது}, \emph{ēṟpaṭṭatu}.

With -\ta{க்கு} for the one affected: \ta{அவருக்குக்} \ta{காயம்} \ta{ஏற்பட்டது}.

\begin{cousinweb}
\textbf{\ta{ஏற்படு}} (\emph{ēṟpaṭu}) is native Dravidian: DEDR 905 lists it under \ta{ஏல்}, \textquotedblleft{}to happen, occur\textquotedblright{}, as \textquotedblleft{}to come into existence\textquotedblright{}.
\end{cousinweb}

\begin{grammarlens}[title={What you've built}]
One more word for telling what happened.
\end{grammarlens}

\subsection*{Your turn}
\begin{itemize}
\raggedright
  \item \emph{Say it:} \emph{ēṟpaṭu}
  \item \emph{Say it:} \emph{ēṟpaṭu}, once more
  \item \emph{Say it:} \emph{cālaiyil periya nerical ēṟpaṭṭatu.} — A big traffic jam arose on the road.
  \item \emph{Recall:} say the Tamil for following this, then the Tamil for to take place, then say \emph{ēṟpaṭu} again
\end{itemize}

\begin{tcolorbox}[breakable,colback=teal!4,colframe=teal!35!black,title={Before you move on}]
What does \ta{ஏற்படு} mean? (\textquotedblleft{}To arise; to come about\textquotedblright{}.) Say it once more.
\end{tcolorbox}
\section[\texorpdfstring{\ta{நிகழ்வு} (nikaḻvu)}{நிகழ்வு (nikaḻvu)}]{\ta{நிகழ்வு} (nikaḻvu) — an event}

\label{lesson:TA-C315-nikalvu}



\begin{quote}
Before the new one: say the Tamil for to take place, then the Tamil for to arise.
\end{quote}

\subsection*{You'll want to know: \ta{நிகழ்வு}}
\textbf{\ta{நிகழ்வு}} — \emph{nikaḻvu} — \textquotedblleft{}an event\textquotedblright{}.

\begin{quote}
\ta{அந்த} \ta{நிகழ்வை} \ta{என்னால்} \ta{மறக்க} \ta{முடியாது}.
\end{quote}

\begin{quote}
\emph{anta nikaḻvai eṉṉāl maṟakka muṭiyātu.}
\end{quote}

\begin{quote}
I cannot forget that event.
\end{quote}

From \ta{நிகழ்}, to happen.

\begin{cousinweb}
\textbf{\ta{நிகழ்வு}} (\emph{nikaḻvu}) is native Dravidian: a noun from \ta{நிகழ்}, \textquotedblleft{}to happen, occur\textquotedblright{} (DEDR 3661).
\end{cousinweb}

\begin{grammarlens}[title={What you've built}]
One more word for telling what happened.
\end{grammarlens}

\subsection*{Your turn}
\begin{itemize}
\raggedright
  \item \emph{Say it:} \emph{nikaḻvu}
  \item \emph{Say it:} \emph{nikaḻvu}, once more
  \item \emph{Say it:} \emph{anta nikaḻvai eṉṉāl maṟakka muṭiyātu.} — I cannot forget that event.
  \item \emph{Recall:} say the Tamil for to take place, then the Tamil for to arise, then say \emph{nikaḻvu} again
\end{itemize}

\begin{tcolorbox}[breakable,colback=teal!4,colframe=teal!35!black,title={Before you move on}]
What does \ta{நிகழ்வு} mean? (\textquotedblleft{}An event\textquotedblright{}.) Say it once more.
\end{tcolorbox}
\section[\texorpdfstring{\ta{தற்போது} (taṟpōtu)}{தற்போது (taṟpōtu)}]{\ta{தற்போது} (taṟpōtu) — at present; now}

\label{lesson:TA-C315-tarpotu}



\begin{quote}
Before the new one: say the Tamil for to arise, then the Tamil for an event.
\end{quote}

\subsection*{You'll want to know: \ta{தற்போது}}
\textbf{\ta{தற்போது}} — \emph{taṟpōtu} — \textquotedblleft{}at present; now\textquotedblright{}.

\begin{quote}
\ta{தற்போது} \ta{அவர்} \ta{வெளிநாட்டில்} \ta{இருக்கிறார்}.
\end{quote}

\begin{quote}
\emph{taṟpōtu avar veḷināṭṭil irukkiṟār.}
\end{quote}

\begin{quote}
At present he is abroad.
\end{quote}

A formal word for \ta{இப்போது}.

\begin{cousinweb}
\textbf{\ta{தற்போது}} (\emph{taṟpōtu}) is built from \ta{தன்}, \textquotedblleft{}its own\textquotedblright{}, and \ta{போது}, \textquotedblleft{}time\textquotedblright{} (DEDR 4559); \ta{ன்} turns to \ta{ற்} before \ta{ப}.
\end{cousinweb}

\begin{grammarlens}[title={What you've built}]
One more word for telling what happened.
\end{grammarlens}

\subsection*{Your turn}
\begin{itemize}
\raggedright
  \item \emph{Say it:} \emph{taṟpōtu}
  \item \emph{Say it:} \emph{taṟpōtu}, once more
  \item \emph{Say it:} \emph{taṟpōtu avar veḷināṭṭil irukkiṟār.} — At present he is abroad.
  \item \emph{Recall:} say the Tamil for to arise, then the Tamil for an event, then say \emph{taṟpōtu} again
\end{itemize}

\begin{tcolorbox}[breakable,colback=teal!4,colframe=teal!35!black,title={Before you move on}]
What does \ta{தற்போது} mean? (\textquotedblleft{}At present; now\textquotedblright{}.) Say it once more.
\end{tcolorbox}
\section[\texorpdfstring{\ta{எப்போதாவது} (eppōtāvatu)}{எப்போதாவது (eppōtāvatu)}]{\ta{எப்போதாவது} (eppōtāvatu) — once in a while; ever}

\label{lesson:TA-C315-eppotavatu}



\begin{quote}
Before the new one: say the Tamil for an event, then the Tamil for at present.
\end{quote}

\subsection*{You'll want to know: \ta{எப்போதாவது}}
\textbf{\ta{எப்போதாவது}} — \emph{eppōtāvatu} — \textquotedblleft{}once in a while; ever\textquotedblright{}.

\begin{quote}
\ta{நான்} \ta{எப்போதாவது} \ta{வானொலியில்} \ta{செய்தி} \ta{கேட்பேன்}.
\end{quote}

\begin{quote}
\emph{nāṉ eppōtāvatu vāṉoliyil ceyti kēṭpēṉ.}
\end{quote}

\begin{quote}
Once in a while I listen to the news on the radio.
\end{quote}

In a question it means ever: \ta{எப்போதாவது} \ta{பார்த்திருக்கிறாயா}?

\begin{cousinweb}
\textbf{\ta{எப்போதாவது}} (\emph{eppōtāvatu}) is built from \ta{எப்போது}, \textquotedblleft{}when\textquotedblright{}, on native \ta{போது}, \textquotedblleft{}time\textquotedblright{} (DEDR 4559), with \ta{ஆவது}, \textquotedblleft{}at least, any\textquotedblright{}.
\end{cousinweb}

\begin{grammarlens}[title={What you've built}]
One more word for telling what happened.
\end{grammarlens}

\subsection*{Your turn}
\begin{itemize}
\raggedright
  \item \emph{Say it:} \emph{eppōtāvatu}
  \item \emph{Say it:} \emph{eppōtāvatu}, once more
  \item \emph{Say it:} \emph{nāṉ eppōtāvatu vāṉoliyil ceyti kēṭpēṉ.} — Once in a while I listen to the news on the radio.
  \item \emph{Recall:} say the Tamil for an event, then the Tamil for at present, then say \emph{eppōtāvatu} again
\end{itemize}

\begin{tcolorbox}[breakable,colback=teal!4,colframe=teal!35!black,title={Before you move on}]
What does \ta{எப்போதாவது} mean? (\textquotedblleft{}Once in a while; ever\textquotedblright{}.) Say it once more.
\end{tcolorbox}
